\documentclass[11pt]{article}
\usepackage[margin=25mm]{geometry}
\usepackage{fontspec,unicode-math}
\setmainfont{Latin Modern Roman}
\setmathfont{Latin Modern Math}
\usepackage{amsmath,amssymb,mathtools}
\usepackage{xurl,hyperref}
\hypersetup{colorlinks=true,urlcolor=blue,linkcolor=blue}
\setcounter{secnumdepth}{0}
\setlength{\parindent}{0pt}
\setlength{\parskip}{5pt}
\setlength{\emergencystretch}{3em}
\newcommand{\sref}[2]{\hyperlink{source:#1:#2}{[S#2]}}
\newcommand{\cataloguescope}[1]{\paragraph{Source question.}#1}
\newcommand{\floor}[1]{\left\lfloor#1\right\rfloor}
\begin{document}
\section{Universal Approximation by One-Lipschitz Spline Networks}\label{statement}

\subsection{Definitions and mathematical statement}

\paragraph{Shared notation.}Standard number systems are denoted by $\mathbb N,\mathbb Z,\mathbb Q,\mathbb R,\mathbb C$; index ranges, computational models and specialized notation follow the cited problem source and the definitions below.


A map $f$ is one-Lipschitz for specified input/output norms if $\lVert f(x)-f(y)\rVert\leq\lVert x-y\rVert$ for all $x,y$. A spline activation is continuous and piecewise affine; a scalar one-Lipschitz spline has every affine slope in $[-1,1]$. The source architecture uses learnable activations with at least three linear regions and constrained affine-layer weights. Universality means every target in the intended one-Lipschitz class can be approximated arbitrarily accurately in the source approximation topology; on a fixed compact domain, uniform accuracy is $\sup_x\lVert f(x)-g(x)\rVert<\varepsilon$. Depth counts successive layers and width counts units in a layer.
\paragraph{Statement recorded in the supplied source.}
Are weight-constrained neural networks with learnable one-Lipschitz spline activations universal approximators of one-Lipschitz functions? Does universality extend to vector-valued targets, and what depth-width structures are optimal? \sref{statement}{2}
\paragraph{Scope qualification.}
The weight constraints and norms are essential architectural data supplied by the cited spline-network paper; the question does not refer to arbitrary unconstrained neural networks. Compact-domain uniform accuracy above explains a standard approximation convention rather than imposing a new global-domain theorem.
\paragraph{Source status.}
The supplied catalogue labels this record open. Its source-status metadata was checked against the supplied snapshot on 2026-10-11. This import does not independently establish current open status or mathematical correctness. The catalogue literature assessment is dated 2026-08-22.
\subsection{Short English statement}
Can the specified Lipschitz-constrained spline networks approximate every one-Lipschitz target, including vector outputs, and how much depth or width is necessary?
\subsection{Sources}
\begin{itemize}
\item[\hypertarget{source:statement:1}{S1}] Ulam AI and the UnsolvedMath Contributors, UnsolvedMath: A Curated Collection of Open Mathematics Problems, catalogue entry. Catalogue curation and metadata: CC BY 4.0. The supplied source question is retained; definitions and the English summary are editorial additions.
\url{https://huggingface.co/datasets/ulamai/UnsolvedMath}.
\item[\hypertarget{source:statement:2}{S2}] Mathematical Imaging and Surface Processing (2023), Oberwolfach Reports, DOI 10.4171/owr/2022/38, pp. 2195--2266
\url{https://doi.org/10.4171/owr/2022/38}.
\item[\hypertarget{source:statement:3}{S3}] Neumayer--Goujon--Bohra--Unser, deep splines
\url{https://doi.org/10.1137/22M1504573}.
\textit{Underlying Oberwolfach Reports material retains its source-specific licence, including CC BY-SA 4.0 where applicable; see THIRD\_PARTY\_NOTICES.md.}
\end{itemize}
\paragraph{Preparation and review.}
Original question attributed to the cited source and UnsolvedMath contributors. Codex (GPT-6; exact serving revision not exposed) assisted with definitions, English summaries, source inspection and formatting. No new human mathematical authorship or independent proof verification is asserted. Catalogue identity checked against the complete supplied registry and existing repository statements; known flagged repetitions and obvious identical conjectures excluded. Full mathematical-equivalence/literature review has not been performed. Statement preparation and final review followed contribution prompts C and D; automated checks verify structure and identity, not open status.
\label{end:statement}
\end{document}
